\chapter{Methodology}
\label{ch:method}

\section{Research design}

The project is a modelling and simulation study whose organising principle is
\term{staged verification}: the work proceeds as a chain from manufacturer
datasheet to switched circuit, and no link in the chain is used until it has
been checked against something that does not share its assumptions.

That principle is a response to the specific hazard of this kind of work. A
simulation study rarely fails by crashing. It fails by producing numbers that
are plausible, internally consistent and wrong, and such failures are invisible
unless each stage is independently validated. Every methodological choice below
follows from taking that hazard seriously.

\begin{table}[htbp]
\centering
\caption{The verification chain: what each stage is checked against.}
\label{tab:verifychain}
\begin{tabular}{@{}L{0.27\textwidth}L{0.34\textwidth}L{0.29\textwidth}@{}}
\toprule
\textbf{Stage} & \textbf{Checked against} & \textbf{Criterion met} \\
\midrule
Parameter identification & the manufacturer datasheet, all five constraints & residual \num{1.4e-14}\,\si{\percent} \\
Analytic $\mathrm{d}I/\mathrm{d}V$ & central differences of the solved current & relative error $<\num{1e-5}$ \\
Maximum power point solver & stationarity of the analytic gradient & $|\mathrm{d}P/\mathrm{d}V| < \num{1e-6}$ \\
Converter linearisation & central differences of the nonlinear vector field & relative error $<\num{1e-9}$ \\
Modal formulas & eigenvalues of the state matrix & agreement to \num{1e-8} \\
Cascade design & margins recomputed on the full-order plant & PM $>\SI{85}{\degree}$ \\
Reference tracking & the analytic I--V curve at the settled point & error $<\SI{0.05}{\percent}$ \\
Averaged converter model & a switched circuit simulation & voltage \SI{0.23}{\percent}, power \SI{0.005}{\percent} \\
Integrator step size & the same run at half the step & final voltage within \SI{0.5}{\volt} \\
\bottomrule
\end{tabular}
\end{table}

\section{Choice of modelling approach}

\subsection{Why a physics-based model rather than a curve fit}

The photovoltaic array could have been represented by a lookup table or a
polynomial fit to published curves. A physics-based single-diode model was
chosen for three reasons. It extrapolates correctly to irradiance and
temperature combinations not present in the datasheet, which a fit cannot. It
yields analytic derivatives, which the control analysis needs and which
differentiating a fitted curve would produce only noisily. And its parameters
are physically interpretable, so an implausible result, such as an ideality
factor of \num{3}, is recognisable as an error rather than accepted as a number.

\subsection{Why the De Soto closure}

Section~\ref{sec:gap} identified the choice between fixing the ideality factor
and solving a determined system. The determined system was chosen because it
uses the temperature coefficient the datasheet already publishes, leaves no
parameter to the modeller's judgement, and makes the ideality factor an output
that can be checked against physical expectation. The cost is a badly scaled
nonlinear solve, addressed in Section~\ref{sec:numerics}.

\subsection{Why both an averaged and a switched converter model}

The two models answer different questions and neither alone is sufficient.

The averaged model is continuous and low-order, so it can be linearised
analytically, swept over an operating envelope in seconds, and simulated for
tens of seconds of closed-loop operation. Frequency-domain analysis is only
possible on it. But it assumes that switching can be replaced by its average,
which is exactly the assumption a control study should not simply accept.

The switched model makes no such assumption, at the cost of being far slower and
not admitting linear analysis. Its role is therefore not to produce the results
but to \emph{audit} them. Using both, and cross-checking, is what allows the
frequency-domain conclusions to be stated with confidence.

\subsection{Why the DC bus is treated as stiff}

In a real string inverter the grid-side stage regulates the DC bus, and its
control bandwidth is far below the converter loops studied here. Modelling the
bus as an ideal voltage sink isolates the MPPT control problem from the grid-side
problem, which is a separate design task. The limitation this imposes is stated
in Chapter~\ref{ch:physical}.

\section{Choice of controlled variable and architecture}

Two decisions were made on analytical grounds rather than convention, and both
are justified where they are made rather than asserted here.

\begin{enumerate}[leftmargin=1.6em]
  \item \textbf{Photovoltaic voltage is the controlled variable}, because
        $\Vmp$ is nearly invariant with irradiance
        (Chapter~\ref{ch:pvmodel}) and because the duty-to-input-voltage
        transfer has no right half plane zero
        (Section~\ref{sec:nozero}), unlike the duty-to-output-voltage transfer.
  \item \textbf{The architecture is cascaded}, because the inner current loop
        sees an operating-point-independent plant while a single voltage loop
        does not (Chapter~\ref{ch:control}). The single-loop design was
        implemented and analysed as well, so that the comparison is quantitative
        rather than asserted.
\end{enumerate}

The design procedure at each layer was loop shaping: choose a crossover
frequency from the physical constraints, place the integral zero below it, and
solve for the proportional gain that makes the loop gain unity at crossover.
Crossovers were chosen by timescale separation, each layer roughly eight times
slower than the one below, with the innermost an order of magnitude below the
switching frequency.

Critically, the gains were derived from \emph{simplified} models, the
high-frequency limit for the inner loop and the ideal-inner-loop limit for the
outer, and then \emph{verified on the exact full-order plant}. Design on a
simplified model is legitimate only if the simplification is afterwards
discharged.

\section{Experimental design for the algorithm comparison}

\subsection{Controlling the confounds}

Chapter~\ref{ch:litreview} noted that published MPPT comparisons are often
confounded by unequal implementations. Four measures were taken so that measured
differences are attributable to the algorithms.

\begin{itemize}[leftmargin=1.4em]
  \item \textbf{A common interface.} All trackers derive from one abstract class
        that owns reference clamping, limit handling and update timing.
        Subclasses supply only the update rule.
  \item \textbf{Identical plant and inner loops.} Every tracker drives the same
        converter with the same cascaded controller and the same gains.
  \item \textbf{Identical sample period.} All hill-climbers run at
        \SI{20}{\milli\second}. Extremum seeking necessarily runs faster,
        because it must resolve a \SI{25}{\hertz} dither; that difference is
        intrinsic to the method and is stated rather than hidden.
  \item \textbf{Identical scoring.} One efficiency routine, one exact
        maximum-power reference, one exclusion window for the start-up
        transient.
\end{itemize}

\subsection{Choice of test signals}

Three signal classes are used, in increasing order of severity.

\begin{description}[leftmargin=0pt, style=unboxed]
\item[Constant irradiance.] Isolates steady-state cost: limit cycle, dither loss
  and bias. Nothing else contributes.
\item[Step change.] The conventional benchmark. Tests recovery from a sudden
  displacement of the optimum, but leaves irradiance constant on both sides of
  the step, so it does not stress the attribution assumption a hill-climber
  makes.
\item[Trapezoidal ramps.] EN 50530-style profiles at four ramp rates. These
  stress exactly the assumption that a step test leaves untouched, and are the
  reason the ranking of the trackers changes between
  Chapter~\ref{ch:mppt} and Chapter~\ref{ch:benchmark}.
\end{description}

Including all three was a deliberate choice: the disagreement between them is
itself a finding, and a study using only one class would have reported a
different and less reliable conclusion.

\subsection{Parameter selection within each algorithm}

Where an algorithm has a free parameter, it was chosen by sweep rather than by
assertion, and the sweep is reported. The perturb and observe step size was
swept over seven values with acquisition time and steady-state efficiency
measured separately, so that the trade-off is exhibited rather than a single
tuned point being presented. The extremum seeking adaptation gain was swept over
six values, revealing the stability limit predicted by averaging theory.

\subsection{The efficiency metric}

Dynamic MPPT efficiency, equation~\eqref{eq:eta}, was adopted because it is the
standard figure of merit and because it integrates rather than samples, so it
captures both lag and ripple in one number. It requires the true available power
at every instant, which is why the exact maximum power point solver of
Chapter~\ref{ch:pvmodel} had to be exact rather than approximate: any error in
the denominator would appear as a spurious difference between trackers.

Worst instantaneous power deficit is reported alongside it, because an
integrated metric can conceal a large but brief shortfall, and
Chapter~\ref{ch:benchmark} contains exactly such a case.

\section{Numerical methods and their justification}
\label{sec:numerics}

\begin{table}[htbp]
\centering
\caption{Numerical choices and the reason for each.}
\label{tab:numchoices}
\begin{tabular}{@{}L{0.30\textwidth}L{0.62\textwidth}@{}}
\toprule
\textbf{Choice} & \textbf{Reason} \\
\midrule
Log parameterisation of $I_0$, $R_s$, $R_{sh}$ & keeps them positive by construction and conditions a system whose unknowns span ten orders of magnitude \\
Residual scaling by $\Isc$ and $|\beta_{Voc}|$ & makes the five residuals comparable so the trust region is meaningful \\
Multi-start over six ideality seeds & removes dependence on the initial guess \\
Damped Newton with backtracking for the diode equation & the residual is monotone and concave, so Newton converges from $I = 0$; backtracking covers excursions beyond $\Voc$ \\
Warm starting from the previous sample & reduces Newton to two or three iterations, which is what makes second-scale closed-loop simulation feasible \\
Bracketed root solve for $\Voc$ and the MPP & guarantees the correct root rather than the nearest one \\
Fixed-step RK4 at $1/(2f_{sw})$ & deterministic sampling instants for the discrete controllers; convergence verified by halving \\
Interpolated $\Pmp(G)$ for scoring & exact solve at every one of \num{500000} samples would be prohibitive; the function is smooth so interpolation error is negligible \\
\bottomrule
\end{tabular}
\end{table}

\section{Reproducibility}

Three practices support reproducibility, described in full in
Chapter~\ref{ch:software}.

\begin{enumerate}[leftmargin=1.6em]
  \item \textbf{A single configuration source.} Array topology, component values
        and crossover frequencies are defined in one function that every script
        consumes, so no figure can silently disagree with another.
  \item \textbf{The Simulink model is built by a script.} A binary model file
        cannot be reviewed or diffed; a builder script can, and it regenerates
        the model from source on demand.
  \item \textbf{One command runs everything.} \texttt{run\_all} regenerates
        every figure and table in this report and runs the test suite. Every
        number quoted in this report came from a single execution of that
        pipeline.
\end{enumerate}

Verification is automated as 27 unit tests. Several encode the report's
conclusions as executable assertions rather than merely checking mechanics: that
the duty-to-power gain vanishes at the maximum power point, that the cascade
holds above \SI{45}{\degree} of phase margin while the single-loop design falls
below \SI{10}{\degree}, and that the perturb and observe limit cycle equals twice
the step. A future change that quietly invalidates an argument in this report
fails a test.

\section{Scope boundaries}

The following were excluded deliberately, and the reasons are given so that the
boundaries are not mistaken for oversights.

\begin{table}[htbp]
\centering
\caption{Deliberate exclusions.}
\label{tab:exclusions}
\begin{tabular}{@{}L{0.30\textwidth}L{0.62\textwidth}@{}}
\toprule
\textbf{Excluded} & \textbf{Reason} \\
\midrule
Partial shading and global search & changes the problem from local tracking to global optimisation; a separate study \\
Sensor noise and quantisation & would strengthen the case for incremental conductance; omitting it is stated as the reason no advantage is observed \\
Grid-side inverter stage & regulates the bus far more slowly than the loops studied; separable by timescale \\
Switching transitions, parasitics, electromagnetic interference & irrelevant to control bandwidth, and modelling them would require device-level data \\
Thermal dynamics of the module & cell temperature varies on a timescale of minutes, far slower than any loop here \\
Hardware construction & outside the time available; specified instead in Chapter~\ref{ch:physical} with quantitative acceptance criteria \\
\bottomrule
\end{tabular}
\end{table}
